\subsection{Causal representation learning and environment
inference}\label{causal-representation-learning-and-environment-inference}

Schölkopf et al.~\cite{scholkopf2021} cast
representation learning as a causal problem and organize the field
around a divide between representations learned under i.i.d. sampling
and those tied to an underlying causal model. Identifiability sits at
the center: without a structural anchor, a representation is determined
only up to a broad class of transformations, and downstream causal
claims inherit the ambiguity.

That anchor has taken several forms. Classical linear ICA establishes
source identifiability up to permutation and scaling under independence
and suitable non-Gaussianity
conditions~\cite{comon1994}. Hyvärinen and
Morioka~\cite{hyvarinen2017} identify a nonlinear
independent-component model from temporal dependence alone.
Interventions supply a further anchor. Ahuja et
al.~\cite{ahuja2023} show interventional data identify
latent variables up to componentwise nonlinearities, Squires et
al.~\cite{squires2023} pin down how many interventions
the linear case requires, and Varıcı et
al.~\cite{varici2025} reach identification under
sparsity and faithfulness through score matching. These results obtain
identification from different sources of structure, including
independence, temporal dependence, and interventions, and they concern
latent representations under their respective observation models. They
do not share a single premise that a known environment index is
required. Our focus is instead on what unlabeled pooled observations
determine about a latent decomposition, and on the additional conditions
required for recording-level grouping.

A related line studies multi-environment learning and asks what happens
when the environment partition is not handed to the analyst. Rosenfeld
et al.~\cite{rosenfeld2021} show that a solved invariant
objective can admit near-optimal solutions that do no better than
empirical risk once the environments shift. Creager et
al.~\cite{creager2021} infer a partition maximally
informative for downstream invariant learning and score it against
worst-group performance. Yao et al.~\cite{yao2025}
formulate multi-environment learning under a known index and an explicit
uniqueness condition, and Murata et
al.~\cite{murata2025} extend invariant-risk objectives
to a clustered setting and characterize the environmental complexity
needed for recovery. Ng et al.~\cite{ng2025} study
causal representation learning from general environments under
nonparametric mixing transformations and state the conditions under
which the components separate. These partitions serve a predictive or
invariance objective. The question we keep separate is whether groups
can be recovered without such an objective defining them, under the two
observation settings made precise below.

\subsection{Mixture identifiability from distributions and
moments}\label{mixture-identifiability-from-distributions-and-moments}

A separate line asks what a single pooled distribution reveals about its
components. A finite mixture of normal distributions is identifiable up
to permutation from the full
distribution~\cite{teicher1963}, and Yakowitz and
Spragins~\cite{yakowitz1968} state the general
mechanism: identifiability follows whenever the component families are
linearly independent. Both identify the component parameters and say
nothing about the assignment of individual observations.

An unlabeled finite-mixture model can also be identified from marginal
independence without observing the partition. Kanamori et
al.~\cite{kanamori2026} recover the latent components
and the mixing matrix from such independence, a route that relies on
different identifying assumptions from the covariance-based construction
in this paper.

Identifiability from the full distribution does not imply
identifiability from a fixed collection of moments. A single pooled
covariance, which is one second moment, is underdetermined. For centered
Gaussian mixtures, the injectivity of moment maps depends on the model
class, the number of components, and the available moment
order~\cite{blomenhofer2023}. The fourth-order map is
non-injective for the generic class through the parallelogram identity,
and injectivity at higher order is established within explicit rank and
parameter-space conditions, not as a universal threshold for every real
positive-definite model. What those results identify is the unordered
collection of component quadratic forms, not a partition or a
finite-sample statement. Our result concerns a restricted fourth-order
inversion in a two-dimensional, equally weighted, zero-mean Gaussian
mixture with component covariances \(I+K\) and \(I-K\). It identifies
the unordered covariance pair and does not establish recovery of
individual labels.

The vocabulary of partial identification is useful when a single
parameter value is no longer determined.
Manski~\cite{manski2003} provides a systematic treatment
in which the target is a set of parameter values compatible with the
observed distribution and the maintained assumptions. Chen et
al.~\cite{chen2026} apply the same set-valued logic to a
latent prevalence and construct bounds from restricted information,
which is the idea we borrow when describing the covariance structures
consistent with relaxed assumptions. We use the term descriptively and
do not claim a sharp set, a set-valued estimator, or a coverage result.

\subsection{Clustering theory and covariance
representations}\label{clustering-theory-and-covariance-representations}

Environment discovery connects the learning-theoretic work above to the
clustering mechanism that performs the recovery. Our unrestricted
partition example illustrates why uniqueness requires explicit
structural assumptions when the grouping is unobserved; it does not
refute the uniqueness assumptions used inside structured causal models.

The statistical foundation is the K-means theory.
Pollard~\cite{pollard1981} proves strong consistency of
the K-means minimizer, so under moment and nondegeneracy conditions the
empirical solution converges almost surely to a population optimum. The
Lloyd iteration~\cite{lloyd1982} is the algorithm in
everyday use, and finite-initialization Lloyd solutions retain an
explicit optimization-error term that is not controlled here. Stability
and approximation results add a complementary picture. Rakhlin and
Caponnetto~\cite{rakhlin2007} analyze the stability of
K-means clustering under distributional perturbations, and Ostrovsky et
al.~\cite{ostrovsky2012} give conditions under which
Lloyd-type methods are effective and the cluster structure is well
separated. Pollard's consistency delivers convergence to a population
optimum but stays silent on whether that optimum aligns with the true
environment, which is the point where our finite-sample argument begins.
We state population alignment as an explicit compatibility premise and
build a risk-gap guarantee on top of it, then convert the excess risk
into a control on assignment error.

The signal on which clustering runs has its own geometry. The matrix
logarithm turns positive-definite matrices into Euclidean coordinates
within the log-Euclidean framework~\cite{arsigny2007}. A
separate construction, the affine-invariant Riemannian
metric~\cite{pennec2006}, gives a different geometry on
the same cone. Our pipeline uses averaged log-covariance features and
Euclidean K-means. Ledoit and Wolf~\cite{ledoit2004}
provide the linear shrinkage estimator used to stabilize the sample
covariance, and the later nonlinear spectral
construction~\cite{ledoit2012} is a separate estimator
we do not adopt.

The empirical setting sits inside a mature EEG literature. Lotte et
al.~\cite{lotte2018} survey EEG-based BCI classification
algorithms, and the PhysioNet resource is documented by Goldberger et
al.~\cite{goldberger2000}. The label we recover is a
source-associated technical structure, not a state of the brain or a BCI
classification result.

Our analysis distinguishes two observation settings. Under pooled
observations, we study unrestricted partition non-identifiability and a
constrained fourth-order identification result for component covariance
parameters. Under repeated recordings, we study how excess K-means risk
controls assignment error under explicit population-alignment and
regularity assumptions. The EEG experiments evaluate source-associated
grouping, cluster-number selection, and sensitivity to the analysis
design. The covariance estimators, feature mappings, and clustering
algorithms are established components used to investigate these
questions.
